% ---------------------------------------------------------------------
%  §5. Conclusiones
% ---------------------------------------------------------------------
\section{Conclusiones}

\begin{enumerate}[label=\arabic*.]
  \item \textbf{Nivel de desarrollo:} el impacto no es igual para todos. Los estudiantes de países en
  desarrollo presentan más horas de uso, menos sueño, peor salud mental y mayor afectación académica percibida.
  \item \textbf{Sueño y rendimiento:} a mayor uso, menos sueño. Tres de cada cuatro estudiantes que se sienten
  afectados académicamente duermen menos de 7 horas, contra uno de cada ocho entre quienes no. Esto es
  coherente con que parte del impacto académico se dé a través de la falta de sueño, aunque el análisis no
  permite confirmarlo.
  \item \textbf{Plataformas y género:} WhatsApp, TikTok y Snapchat se asocian con peores resultados, y
  LinkedIn, Twitter y Facebook con resultados más favorables, pero gran parte de la diferencia se explica por
  las horas de uso. La brecha de género es muy leve.
  \item \textbf{Relaciones sentimentales:} el uso diario es prácticamente igual entre solteros, en pareja y en
  situación complicada; las diferencias en la plataforma elegida describen el perfil de sus usuarios, no un
  efecto de la red sobre las relaciones.
\end{enumerate}

\begin{hallazgo}
\textbf{Conclusión general.} Las cuatro preguntas convergen en un mismo patrón: \textbf{la cantidad de horas
de uso es el factor que más se repite detrás de los resultados menos favorables}. Los grupos con peores
indicadores son también los que más horas pasan en las redes, y a más horas se observan menos sueño y menor
salud mental. Más que \emph{qué} red se usa o \emph{quién} la usa, el dato más útil para anticipar un impacto
negativo es \textbf{cuánto tiempo se le dedica}.
\end{hallazgo}

\textbf{Limitaciones.} Los resultados son descriptivos y no permiten establecer causalidad; el rendimiento
académico es autodeclarado y no una medición de las notas; la clasificación por desarrollo se basa solo en el
ingreso y los grupos están desbalanceados (1.103 contra 638); varias plataformas tienen pocos casos y el estado
sentimental falta en el 59,55\,\% de los registros; y las preguntas 1 y 2 incluyen los 705 registros
duplicados, aunque sin ellos las conclusiones se mantienen.

\textbf{Extensiones.} Para confirmar la relación haría falta un estudio que siga a los mismos estudiantes a lo
largo del tiempo y mida el rendimiento académico con datos objetivos, en lugar de la percepción de cada
estudiante.
